\section{Conclusion}
\label{sec:conclusion}

This pre-registered survey experiment provides the first large-scale experimental evidence on public perceptions of algorithmic redistricting. Three findings stand out. First, algorithmic maps are perceived as substantially fairer than enacted maps by a broad majority of the American public---a $0.41$ SD advantage that holds across partisan groups and persists after controlling for pre-treatment attitudes. Second, transparency amplifies this advantage: a plain-language description of the algorithmic process adds $0.22$ SD to perceived fairness, a finding with direct implications for how reform advocates should communicate about algorithmic redistricting. Third, there is no significant difference in public evaluations of algorithmic and commission-drawn maps, suggesting that algorithmic redistricting has already cleared the legitimacy threshold established by the independent commission model.

These results matter for redistricting politics for two reasons. The first is tactical: reform advocates seeking to build coalitions behind algorithmic redistricting can credibly claim bipartisan public support. Our results show positive evaluations from strong Democrats, strong Republicans, and independents alike. The partisan gap in evaluations is real but small ($0.18$ SD), and even the group with the smallest effect (strong Republicans) endorses algorithmic maps over enacted maps by a substantial margin.

The second reason is normative. Redistricting reform faces a legitimacy problem that is partly self-created: because gerrymandering benefits one of the two major parties at any given time, the party in power has incentives to describe reform efforts as partisan attacks on its legitimate electoral advantages. Algorithmic redistricting is structurally resistant to this characterization because it is equally likely to cost any party advantage it currently holds. Our results suggest that ordinary citizens perceive this neutrality, even without detailed technical knowledge. When the public can identify algorithmic neutrality and rate it as fair, algorithmic redistricting gains a political resource that purely technical arguments cannot provide.

Reform advocates should draw two concrete lessons from our results. First, \textit{invest in communication infrastructure}: our transparency treatment worked because it was written in plain language and emphasized procedural guarantees rather than technical sophistication. Organizations working on algorithmic redistricting should develop standardized plain-language explanations for public communication that emphasize what the algorithm cannot do and what it guarantees. Second, \textit{pair algorithmic methods with existing commission structures}: our finding that public evaluations of algorithmic and commission-drawn maps are equivalent suggests that the most politically viable path forward may be algorithmic redistricting implemented \textit{by} independent commissions---combining the institutional legitimacy of the commission model with the structural neutrality of algorithmic line-drawing.

Future research should examine whether these perceptions persist under adversarial conditions, including elite cueing by partisan actors who oppose reform, and whether they generalize to respondents who know they personally live in the districts being evaluated. Longitudinal panel studies tracking perception changes as algorithmic redistricting moves from academic proposal to legislative debate would be especially valuable. The present study establishes that the public legitimacy foundation for algorithmic redistricting exists; further research should explore its durability under political pressure.
